% ── ABSTRACT ─────────────────────────────────────────────────────────
\chapter*{Abstract}
\addcontentsline{toc}{chapter}{Abstract}

The rapid evolution of software engineering education demands innovative tools that bridge the gap between isolated individual study and real-world collaborative development practices. This project presents \textbf{AlgoRiddle}, a full-stack, real-time collaborative coding platform designed to transform the way students practice Data Structures and Algorithms (DSA). Unlike conventional online judges that offer a solitary experience, AlgoRiddle enables multiple users to join a shared virtual room where they can simultaneously write and edit code in a live-synchronized editor, brainstorm solutions on a shared multi-tool whiteboard, and communicate through peer-to-peer voice chat---all within a single, unified web interface.

The platform is architected using a modern technology stack comprising React 19 with Vite on the frontend, Node.js with Express on the backend, MongoDB Atlas for persistent data storage, Socket.IO for bi-directional real-time event communication, and native WebRTC for browser-based voice chat. Code execution is handled through integration with the Wandbox compilation API, supporting JavaScript, Python, and C++ with automated test case evaluation. The system features a gamified question progression model, story-driven DSA challenges, and an intuitive dark-themed user interface that delivers a premium development experience.

This report details the complete software development lifecycle of AlgoRiddle, encompassing the problem formulation, background research, literature survey, system design, implementation methodology, and comprehensive testing results. The platform has been successfully deployed using a zero-cost cloud infrastructure stack (Render, Vercel, MongoDB Atlas) and demonstrates the feasibility of building production-grade collaborative educational tools using modern web technologies.

\clearpage
